\documentclass[conference]{IEEEtran}
\usepackage{amsmath,amssymb}
\usepackage{booktabs}
\usepackage{array}
\usepackage{url}

\title{Can Quantum Artificial Intelligence Imitate the Human Brain?\\
A Design-Only Benchmark for Resource-Matched Neural Modeling}

\author{\IEEEauthorblockN{Four-Agent Research Workflow}
\IEEEauthorblockA{Survey -- Idea -- ExperimentDesign -- Author\\
Design-only research report}}

\begin{document}
\maketitle

\begin{abstract}
The human brain is a multiscale biological system whose neurons, synapses, dynamics,
and learning rules are only partially understood. Quantum artificial intelligence may
offer new computational representations, but the existence of a quantum algorithm is
not evidence that the brain uses quantum computation. This paper proposes
Q-BRAIN-IMITATION-BENCHMARK, a design-only research program that separates quantum
machine learning, brain-like artificial intelligence, and quantum-brain hypotheses. It
compares quantum kernels, variational circuits, hybrid quantum-classical models,
classical recurrent networks, and spiking networks on declared targets such as recorded
neural populations, validated microcircuit simulators, and closed-loop neural-style
agents. Evaluation covers signal fidelity, dynamical fidelity, behavior, learning,
held-out transfer, and practical quantum resources. Ideal simulation, noisy simulation,
and hardware-aware execution are not pooled. Entanglement, encoding, parameter count,
data access, and optimization are separately ablated. A quantum advantage is accepted
only when it survives matched classical and spiking controls and a realistic resource
account. A brain-imitation claim additionally requires multi-axis dynamical and transfer
evidence. The protocol makes no claim about consciousness, personhood, or whole-brain
equivalence, and reports no empirical outcomes.
\end{abstract}

\begin{IEEEkeywords}
quantum artificial intelligence, quantum machine learning, brain simulation, spiking
neural networks, neural dynamics, quantum advantage, reproducible evaluation
\end{IEEEkeywords}

\section{Introduction}
\label{sec:introduction}

The question ``Can quantum artificial intelligence imitate the human brain?'' contains
three ideas that are often merged in popular accounts. Quantum machine learning (QML)
uses quantum hardware or quantum circuits to learn patterns and solve computational
problems. Brain-like artificial intelligence models neural representations, spikes,
memory, behavior, or learning. Quantum-brain hypotheses propose that biologically
relevant quantum coherence, collapse, or geometry is necessary for cognition or
consciousness. These are not interchangeable claims.

QML is an active research area. Biamonte \emph{et al.} describe quantum algorithms and
software primitives that may support machine learning, while also emphasizing substantial
hardware and software challenges \cite{biamonte2017}. This motivates a careful question
about whether quantum representations are useful for neural data. It does not establish
that a quantum computer is a brain model or that it will be faster on a practical neural
task.

The biological claim is stronger. A model can predict a recorded spike train while
using internal states unrelated to membrane potentials or synaptic plasticity. A quantum
circuit can approximate a probability distribution without reproducing the causal
mechanisms that created it. At the other extreme, a biophysically detailed simulator
may reproduce local dynamics while failing at behavior or learning. An answer must
therefore declare the target scale and the level of imitation.

The most controversial route is to infer quantum cognition from quantum effects in the
brain. Tegmark calculated short decoherence time scales for degrees of freedom proposed
to be relevant to cognitive processes and argued that a classical treatment of neural
network dynamics is not thereby invalid \cite{tegmark2000}. Hameroff and Penrose review
the Orch OR proposal, which assigns a possible role to quantum processes in microtubules
and spacetime structure \cite{hameroff2014}. The disagreement is scientifically useful
only if each proposal yields a measurable prediction that is compared with classical
stochastic and spiking alternatives.

This paper reframes the question as follows:

\begin{quote}
Under matched task, data, energy, and latency budgets, can a quantum or hybrid model
reproduce the signal, dynamics, learning, and transfer properties of a specified neural
target better than strong classical and spiking baselines, and is any gain attributable
to quantum resources rather than model size, data, or optimization?
\end{quote}

The contributions are fourfold. First, we provide a survey that separates QML, neural
simulation, and quantum-brain evidence. Second, we propose Q-BRAIN-IMITATION-BENCHMARK,
a multi-level evaluation architecture. Third, we specify formal metrics and an ablation
protocol for quantum resources, noise, and practical cost. Fourth, we define a
decision matrix that limits conclusions to the observable target and prevents functional
performance from being interpreted as consciousness. This paper is DESIGN\_ONLY: no
model has been run and no human-derived result is asserted.

\section{Survey of the Evidence}
\label{sec:survey}

\subsection{Quantum machine learning}

The central QML promise is not that a quantum state is biologically similar to a brain.
It is that quantum feature maps, kernels, sampling procedures, or variational circuits
may represent or manipulate information in ways that are difficult for a classical
algorithm to reproduce efficiently. The Nature review by Biamonte \emph{et al.}
organizes this promise around quantum algorithms and applications, while noting that
hardware and software challenges remain considerable \cite{biamonte2017}. The relevant
experimental variables are therefore qubit count, circuit depth, entangling-gate count,
measurement shots, noise, compilation, and classical overhead.

An ideal circuit result is only one point in that resource space. A practical comparison
must include a noise model, finite measurement shots, data movement, circuit compilation,
and the classical optimizer. If a classical neural network receives many more parameters
or a stronger optimizer, a higher score cannot be attributed to quantum structure. The
survey therefore treats resource matching as part of the scientific claim rather than
as an implementation detail.

\subsection{Neural modeling across scales}

Brain imitation has no single ground truth. At the signal level, a model may predict
spike counts, local field potentials, or population activity. At the dynamical level,
it may reproduce synchrony, attractor transitions, adaptation, and perturbation
recovery. At the behavioral level, it may select actions or recall sequences. At the
learning level, it may show retention, forgetting, few-shot adaptation, and transfer.
These levels can disagree.

Large-scale reconstruction work demonstrates why scale must be explicit. Markram
\emph{et al.} reconstruct and simulate a neocortical microcircuit using anatomical and
physiological constraints \cite{markram2015}. Such work does not imply that a complete
brain has been reconstructed. It does show that mechanistic claims require declared
connectivity, parameters, observation models, and perturbations. A benchmark should
therefore compare a quantum model with classical and spiking models on a target that is
small enough to audit but rich enough to expose dynamics.

\subsection{Quantum-brain proposals}

The quantum-brain question is not answered by adding qubits to a neural network.
Tegmark's decoherence analysis argues that cognitive degrees of freedom should generally
be treated as classical on relevant time scales, with decoherence much faster than neural
dynamics under the assumptions studied \cite{tegmark2000}. Hameroff and Penrose's review
offers a different theoretical account involving orchestrated objective reduction and
microtubules \cite{hameroff2014}. The existence of such a proposal is evidence that the
hypothesis is discussable, not evidence that it is true.

The correct experimental consequence is demanding: a quantum-brain theory must specify
an observable, a temporal and spatial scale, an intervention, and a direction of effect.
The observation must be distinguished from classical stochastic noise, ordinary quantum
chemistry that does not support brain-scale computation, and spiking-network dynamics.
Without this comparison, a quantum model that fits data is merely another flexible model.

\subsection{Generalization and the imitation boundary}

Lake \emph{et al.} argue that strong benchmark performance does not automatically give a
system human-like compositionality, causal learning, intuitive physics, or few-shot
learning \cite{lake2017}. These concerns apply directly here. A model trained on one
stimulus distribution may interpolate neural activity while failing a new stimulus
family or perturbation. Held-out tasks and controlled rule changes are needed to
separate fit from a reusable learning process.

The evidence thus supports a bounded research program: QML may be useful for specified
neural tasks; classical and spiking models remain indispensable controls; and quantum
effects in the brain remain a hypothesis requiring biological prediction. It does not
support claims of quantum consciousness, digital personal identity, or whole-brain
equivalence.

\begin{table}[t]
\caption{Claims and evidence boundaries}
\label{tab:claims}
\centering
\footnotesize
\begin{tabular}{p{0.29\columnwidth}p{0.60\columnwidth}}
\hline
Claim & Required evidence in this study \\
\hline
QML utility & Improvement on a declared neural task over matched classical and spiking models. \\
Quantum-resource effect & Degradation after resource ablations and persistence under realistic noise and overhead. \\
Brain-like imitation & Multi-axis match of signal, dynamics, behavior, learning, and held-out transfer. \\
Quantum brain & A biological observable predicted by a quantum model and not explained by matched alternatives. \\
Consciousness & Outside the benchmark; never inferred from task performance. \\
\hline
\end{tabular}
\end{table}

\section{Idea: Q-BRAIN-IMITATION-BENCHMARK}
\label{sec:idea}

The proposed benchmark has a common input encoder, a model-specific computational core,
and a common readout. The encoder maps stimuli or upstream activity to a declared
representation. The core is one of a quantum kernel, variational circuit, hybrid
quantum-classical recurrent model, parameter-matched classical recurrent model, or
spiking neural network. The readout maps states to neural observations, actions, or
memory reports. All conditions receive the same examples, feedback, task instances,
and stopping rule.

The benchmark has four evaluation layers. The signal layer asks whether observed neural
activity can be predicted. The dynamical layer asks whether multi-step trajectories,
synchrony, attractors, and perturbation recovery are reproduced. The behavioral layer
asks whether a task policy or memory behavior is matched. The learning layer asks whether
the model adapts, retains, forgets, and transfers under new inputs. A resource layer is
reported alongside them rather than hidden in a score.

Quantum resources are explicit. For each circuit, record qubits, depth, entangling gates,
shots, noise channel, compilation time, measurement time, classical optimizer steps,
and wall-clock latency. The first ablation replaces the quantum feature map with a
matched classical feature map. The second removes entangling operations. The third
shuffles the quantum encoding while preserving shape and parameter count. These tests
ask whether an apparent gain comes from entanglement, nonlinear representation, model
capacity, or data processing.

The target is not ``the brain'' as a single object. Candidate target families are:

* a recorded neural population with a public stimulus-response split;
* a validated spiking microcircuit with known connectivity and interventions;
* a sequence-memory circuit with controlled noise and delays;
* a closed-loop neural-style controller for a simple, safe digital agent;
* a quantum-sensitive biological probe only when a theory supplies an observable
  prediction.

The main hypotheses are:

* H0: apparent quantum advantage disappears under matched classical surrogates and
  resource controls;
* H1: a hybrid model improves at least one declared neural metric under matched data and
  compute budgets;
* H2: quantum resources affect representation or sample efficiency more readily than
  biophysical fidelity;
* H3: ideal-circuit gains shrink under finite shots, noise, compilation, and latency;
* H4: behavioral similarity can occur without dynamical similarity;
* H5: a quantum-brain claim is supported only when its predicted observable survives
  classical stochastic and spiking alternatives.

\section{Formal Model and Metrics}
\label{sec:formal}

Let $z_{1:T}$ be a target neural trajectory, $\hat{z}_{1:T}$ a model trajectory,
$s_{1:T}$ the input sequence, and $\theta$ the model parameters. A declared observation
model $g$ maps internal states to observed signals. Signal error is represented as
\begin{equation}
E_{\mathrm{sig}}=\frac{1}{T}\sum_{t=1}^{T}d\left(z_t,g(\hat{h}_t;\theta)\right),
\label{eq:signal}
\end{equation}
where $d$ is a predeclared distance and $\hat{h}_t$ is the model state. Likelihood,
correlation, timing error, and calibration are reported as separate choices of $d$ or
complementary metrics; they are not silently averaged.

For spike trains, let $b$ index neurons and $w$ index time bins. A population rate
profile is $r_{bw}$, and the model profile is $\hat{r}_{bw}$. A normalized rate error is
\begin{equation}
E_{\mathrm{rate}}=\frac{\sum_{b,w}|r_{bw}-\hat{r}_{bw}|}{\sum_{b,w}|r_{bw}|+\varepsilon},
\label{eq:rate}
\end{equation}
where $\varepsilon$ is a fixed numerical constant. Sensitivity is reported for at
least two bin widths because a quantum model could appear accurate after temporal
coarsening.

Dynamic fidelity includes autocorrelation, cross-correlation, synchrony, attractor
occupancy, transition probability, and recovery after a controlled perturbation. Let
$q_{uv}$ be the target transition probability between states $u$ and $v$, and let
$\hat{q}_{uv}$ be the model estimate. A transition divergence is
\begin{equation}
E_{\mathrm{trans}}=\frac{1}{2}\sum_{u,v}\left|q_{uv}-\hat{q}_{uv}\right|,
\label{eq:transition}
\end{equation}
which is interpretable as a total-variation distance. The perturbation response is
evaluated on the same input intervention for every model.

Behavioral fidelity is measured by task reward, action agreement, error distribution,
and latency. Learning fidelity includes retention after a delay, forgetting after
interference, and adaptation after a rule change. If $P_{\mathrm{dev}}$ and
$P_{\mathrm{hold}}$ are matched task performances, held-out transfer is
\begin{equation}
T_{\mathrm{hold}}=P_{\mathrm{hold}}-P_{\mathrm{dev}}.
\label{eq:transfer}
\end{equation}
A negative value is informative: it identifies a failure to transfer rather than an
invalid experiment.

To quantify a claimed quantum contribution, let $M_Q$ be the quantum or hybrid model,
$M_C$ its matched classical surrogate, and $M_A$ its resource ablation. For a metric
$P$, define the quantum-resource contrast as
\begin{equation}
\Delta_Q=P(M_Q)-\frac{P(M_C)+P(M_A)}{2}.
\label{eq:quantumcontrast}
\end{equation}
This contrast is descriptive, not a proof of a quantum advantage. Its interpretation
requires matched data, parameters, optimization, and practical cost, together with
uncertainty intervals.

The resource vector is
\begin{equation}
\mathcal{R}=(n_q,L,e,s,\tau_{\mathrm{compile}},\tau_{\mathrm{run}},E),
\label{eq:resources}
\end{equation}
where $n_q$ is qubit count, $L$ circuit depth, $e$ entangling-gate count, $s$ shots,
$\tau$ terms are compilation and execution times, and $E$ is optional energy metadata.
Two systems are resource matched only after the comparison specifies which coordinates
are fixed and which are intentionally varied.

\section{ExperimentDesign}
\label{sec:design}

\subsection{Stage A: controlled offline benchmark}

The first stage uses public neural data or synthetic targets with known ground truth.
Each task card declares target scale, input modality, temporal resolution, observation
model, preprocessing, train/development/held-out split, intervention set, and value
metrics. Held-out data are not merely held-out trials: the benchmark holds out stimulus
families, task rules, intervention magnitudes, and sequence compositions when the target
permits it.

The model conditions are C0, a parameter- and compute-matched classical recurrent
network; C1, a spiking neural network with the same input and readout; Q0, an ideal
quantum kernel or variational circuit; Q1, an ideal hybrid quantum-classical recurrent
model; and Q2, Q1 with finite shots, calibrated noise, compilation, and latency. A1
removes entangling gates, A2 replaces the quantum feature map with a matched classical
map, and A3 shuffles quantum encoding. Hyperparameters are tuned on development tasks
only. Models, compiler versions, seeds, shot counts, optimizer steps, and stopping rules
are frozen before held-out evaluation.

\subsection{Stage B: sandbox interaction}

If interaction is needed, use a local sandbox with a simple digital agent. The human may
provide a task goal or correct a model output, but cannot publish, deploy, purchase, or
make a clinical decision through the interface. Human input is logged as an operation,
and a no-human condition receives the same task description and time budget. The purpose
is to measure learning from feedback, not to infer intention from a conversational style.

\subsection{Stage C: quantum-sensitive biological probe}

This stage is optional and occurs only after a theory specifies a testable observable.
The theory must state the predicted signal, time scale, spatial scale, perturbation, and
expected direction. The analysis includes at least a classical stochastic model and a
spiking model with matched connectivity and noise. Ordinary quantum chemistry is not
automatically evidence for brain-scale quantum computation. Any neural data remain
de-identified and are not used for diagnosis, identity inference, or consciousness
classification.

\begin{table}[t]
\caption{Experimental conditions and purpose}
\label{tab:conditions}
\centering
\footnotesize
\begin{tabular}{p{0.18\columnwidth}p{0.70\columnwidth}}
\hline
Condition & Purpose \\
\hline
C0 & Classical recurrent baseline with matched parameters, data, and optimizer budget. \\
C1 & Spiking baseline with matched input, readout, and training examples. \\
Q0 & Ideal quantum kernel or variational circuit; theoretical upper-bound condition. \\
Q1 & Ideal hybrid quantum-classical recurrent model. \\
Q2 & Finite-shot, noisy, compiled, hardware-aware version of Q1. \\
A1 & Entanglement ablation using separable operations. \\
A2 & Matched classical feature-map surrogate. \\
A3 & Shuffled quantum encoding with preserved dimensions. \\
\hline
\end{tabular}
\end{table}

\subsection{Task protocols}

For neural encoding, the model predicts a held-out population response from a sensory
stimulus. The evaluation includes per-channel likelihood, timing, calibration, and
noise robustness. For dynamical reconstruction, the model receives an initial state
and input sequence and must reproduce multi-step spike statistics, synchrony, and
recovery after an input or state perturbation. For sequence memory, the model learns a
temporal pattern, retains it across a delay, and recalls it after distractor noise.

For closed-loop behavior, a digital agent uses the model output to act in a simple
environment with a fixed safety envelope. The same policy budget applies to all model
families. For few-shot transfer, one stimulus family or task rule is used for
development, while a new family or rule is held out. The number of examples and the
adaptation steps are fixed. A quantum model that wins only after receiving more examples
or more optimizer iterations is not considered resource matched.

\subsection{Randomization and sampling}

Use blocked randomization by task, target scale, noise level, and seed. Generate at least
20 independent runs or candidate trajectories per condition and task instance for the
offline stage. This is an operational minimum for detecting failure modes, not a generic
power guarantee. A development pilot estimates variance and the smallest practically
important effect before the held-out sample is finalized. Ideal and noisy conditions
are analyzed separately.

\section{Ablations, Analysis, and Reproducibility}
\label{sec:analysis}

\subsection{Causal ablations}

The ablation sequence removes one putative source of performance while holding task,
data, model size, optimizer, and seed schedule fixed. Removing entanglement tests one
quantum resource. Replacing the quantum feature map tests whether a classical feature
map can produce the same representation. Shuffling the encoding tests whether the
advantage depends on meaningful input geometry. Q2 tests whether an ideal result remains
after measurement and hardware constraints. A model-selection budget is also fixed:
choosing the best of many circuits after seeing held-out outcomes is prohibited.

\begin{table}[t]
\caption{Ablation questions}
\label{tab:ablations}
\centering
\footnotesize
\begin{tabular}{p{0.27\columnwidth}p{0.62\columnwidth}}
\hline
Removed or changed factor & Diagnostic question \\
\hline
Entangling gates & Does the effect require nonseparable quantum operations? \\
Quantum feature map & Can a matched classical map explain the effect? \\
Input encoding & Is performance tied to a meaningful representation or arbitrary geometry? \\
Ideal to noisy circuit & Does the result survive finite shots and decoherence? \\
Held-out rule & Does the learned behavior transfer beyond the fitted distribution? \\
Dynamic probe & Does behavioral fit conceal incorrect neural dynamics? \\
\hline
\end{tabular}
\end{table}

\subsection{Hierarchical statistical model}

For outcome $y_{ijkt}$ from condition $i$, target family $j$, instance $k$, and seed
$t$, use
\begin{equation}
y_{ijkt}=\mu+\alpha_i+\beta_j+\gamma_k+\delta_t+\eta_{ijk}+\epsilon_{ijkt},
\label{eq:hierarchical}
\end{equation}
where $\alpha_i$ is the condition effect, $\beta_j$ the target-family effect,
$\gamma_k$ the instance effect, $\delta_t$ the seed effect, and $\eta_{ijk}$ an
optional interaction. Report effect sizes, uncertainty intervals, per-task distributions,
and validity rates. A Bayesian implementation declares priors before held-out data are
inspected; a frequentist implementation uses cluster-robust intervals by task instance.

The confirmatory contrasts are Q1 versus C0, Q1 versus C1, Q1 versus A1, Q1 versus A2,
and Q1 versus Q2. The primary metric families are signal fidelity, dynamic fidelity,
behavior, learning, transfer, and resources. Multiplicity is controlled across the
predeclared axes. An exploratory composite profile may be shown only after all axes are
reported and never replaces an axis-level failure.

\subsection{Failure and missing-data policy}

Timeout, invalid output, compiler failure, and evaluator refusal are separate categories.
They are not silently removed. If a trace is missing, the artifact may retain a product
score, but no process claim is made. If a human-derived neural record cannot be shared,
the public package contains schema, hashes, a synthetic surrogate, and the deletion or
access procedure. If a circuit cannot run on hardware, the result is labeled an ideal or
noisy simulation, not hardware evidence.

\subsection{Reproducibility package}

The package releases task cards, public data and license notes, model adapters, quantum
circuits, compiler versions, seeds, optimizer traces, preprocessing, raw event logs,
metric code, exclusion rules, and analysis code. Each run includes model identifier,
input hash, output hash, circuit hash, noise model, shot count, latency, and energy
metadata where available. A manifest maps every reported number to its input hashes and
code version. Reproducibility includes negative and invalid runs rather than only the
best circuit.

\section{Resource Accounting and Practical Execution}
\label{sec:resources}

Quantum neural modeling is often reported at the level of circuit depth or qubit count,
but a brain-like task also pays for encoding, measurement, optimization, and data
movement. A fair comparison includes the classical time needed to construct a circuit,
transfer a neural sequence into it, execute repeated shots, aggregate the readout, and
update parameters. The practical time budget for run $r$ is written as
\begin{equation}
\tau_r=\tau_{\mathrm{encode},r}+s_r(\tau_{\mathrm{queue},r}+\tau_{\mathrm{circuit},r}+\tau_{\mathrm{readout},r})+\tau_{\mathrm{opt},r},
\label{eq:wallclock}
\end{equation}
where $s_r$ is the number of shots and the last term includes the classical optimizer.
The terms are reported separately because a low-depth circuit may be dominated by
repeated measurement or by data encoding.

The benchmark defines three resource regimes. In the ideal regime, Q0 and Q1 use exact
or high-precision simulation and are interpreted as algorithmic feasibility results. In
the noisy regime, a calibrated channel model injects gate, readout, and thermal errors;
shot counts are finite. In the hardware-aware regime, compilation, queueing, device
topology, connectivity, and classical host latency are recorded. A result is not
promoted from ideal to hardware-aware merely because the same circuit can theoretically
be transpiled.

Energy reporting is similarly conditional. If a device exposes power and runtime,
record the measurement method, sampling interval, queue policy, and whether the host
computer is included. If it does not, mark energy as unavailable rather than estimate a
precise value from a nominal chip specification. Classical baselines receive the same
accounting standard, including preprocessing and hyperparameter search. Training and
inference budgets are reported separately.

The input encoding is a central design choice. A population response can be encoded as
amplitudes, angles, basis states, or a classical feature vector followed by a quantum
map. These encodings have different normalization, width, and data-loading costs. The
task card must specify the encoding and include at least one alternative encoding in a
development ablation. If performance changes sharply, the finding is an encoding
result, not automatically an intrinsic property of quantum computation.

The readout also creates a possible confound. A quantum core may be followed by a large
classical recurrent network that performs most of the work. To prevent this, the readout
family, parameter count, optimizer, and recurrent depth are matched across Q1, C0, and
A2. A second analysis reports core-only performance and end-to-end performance. The
difference estimates how much of the result is attributable to the shared classical
head.

Noise can affect neural metrics unevenly. Readout noise may increase spike timing error,
while a coarse population-rate metric remains stable. The benchmark therefore reports
noise sensitivity by metric and temporal resolution. It also reports the number of
repetitions needed to reach a fixed confidence width. A quantum model that needs many
more shots to match a classical rate estimate has a different practical profile from a
model that matches it with fewer shots, even if their final point estimates are equal.

\section{Threats to Validity and Extensions}
\label{sec:threats}

The first threat is target misspecification. A recorded neural population is an
observation through a measurement apparatus, not the complete hidden state of the brain.
A model can match the recording while using an unrelated latent mechanism. The protocol
limits the claim to the declared observation model and adds perturbation response when
the target supports interventions. It does not call a correlation a mechanistic match.

The second threat is simulator circularity. A spiking target generated by a particular
classical simulator can favor another model that shares its inductive assumptions. This
is why the benchmark combines recorded data, multiple simulator families, and
cross-target evaluation. Results are stratified by target source. A model that wins only
on a simulator constructed from its own feature assumptions is not evidence for general
brain imitation.

The third threat is an unfair classical baseline. Classical recurrent models can be
under-tuned, while quantum models may receive an extensive architecture search. The
protocol fixes a search budget, records every tried configuration, and uses the same
development data for all model families. A strong baseline is a scientific requirement,
not an obstacle to a quantum result. If the baseline is updated after seeing held-out
data, the held-out result is invalidated.

The fourth threat is representation dependence. Quantum feature maps and spike-based
representations can make the same input geometry look different. The benchmark publishes
the encoder, tests alternative encoders on development data, and reports sensitivity.
The classical feature-map ablation is especially important: a quantum circuit may offer
a useful nonlinear map even when entanglement is not the cause. That is still a useful
engineering result, but it is a narrower claim.

The fifth threat is multiple testing. Many target scales, metrics, seeds, encodings,
and noise settings create opportunities to select a favorable result. The primary
contrasts and thresholds are declared before held-out evaluation. Exploratory results
are labeled, and all attempted conditions remain in the event log. The benchmark does
not use the best metric from each family to create a post hoc quantum score.

The sixth threat is anthropomorphic interpretation. Words such as memory, attention,
and decision are useful names for task functions, but they do not establish subjective
experience. The report uses ``memory fidelity'' for delayed sequence performance and
``behavioral fidelity'' for action agreement. It does not say that the model remembers,
feels, or is conscious. This linguistic rule is part of the scientific boundary.

Three extensions are suitable after the offline benchmark. A longitudinal extension can
test whether a hybrid model learns a reusable update rule across sessions; checkpoints
and data order must be preserved. An embodied extension can move the controller into a
digital twin and then a supervised laboratory sandbox with an emergency stop. A
multimodal extension can combine spikes, local field potentials, and behavior, but must
declare missingness and alignment. None of these extensions changes the requirement for
matched classical and spiking alternatives.

\section{Measurement Governance and Decision Protocol}
\label{sec:governance}

The benchmark begins with a task card rather than a model. The card identifies the
scientific target, the neural observation, the scale of the claim, the input and output
spaces, the train/development/held-out split, the intervention set, and the tolerances
for each metric. It also identifies which data are public, synthetic, restricted, or
derived. This ordering prevents a convenient dataset or a successful circuit from
defining the research question after the fact.

Each task card has a pre-execution review of five items. First, the target must have a
well-defined observation model. Second, the classical and spiking baselines must be
capable of receiving the same information. Third, the quantum encoding and readout must
be specified, including data-loading cost. Fourth, the held-out split must test a new
concept, rule, or perturbation rather than only a new random seed. Fifth, the primary
contrast and practical threshold must be frozen before held-out execution. A task that
fails one item can remain exploratory but cannot produce a confirmatory quantum claim.

The run manifest is append-only. It includes task identifier, condition, seed, model
version, circuit version, compiler version, input hash, output hash, data split, number
of shots, noise model, optimizer steps, latency, and error class. It stores separate
timestamps for encoding, queueing, execution, readout, and optimization. For a hardware
run, it also stores device topology and calibration date. For an ideal simulation, it
stores simulator precision and explicitly labels the result as ideal. This distinction
protects readers from mistaking an algorithmic demonstration for a device measurement.

Analysis is performed in two passes. The first pass is blind to condition names and
produces metric tables, validity flags, and failure categories. The second pass attaches
the condition labels and estimates the preregistered contrasts. The held-out partition
is not used to choose an encoder, depth, seed, or visualization. Any post hoc analysis
is placed in a clearly labeled exploratory section. The report includes attempted
conditions, not only completed conditions, because circuit compilation failure or
unstable optimization may be a practical property of the model.

The data governance protocol is proportional to the target. Synthetic neural data are
preferred for early circuit experiments because their latent state and interventions
are known. Public recordings are accompanied by license and consent notes. Restricted
recordings are never uploaded to an external model or placed in a public repository;
the release uses a schema, a data dictionary, hashes, and a synthetic surrogate. No
record is used for diagnosis, identity inference, or classification of consciousness.
The benchmark operator records whether any data preprocessing used labels that were
unavailable to the model at deployment time.

The decision report has a fixed order. It first reports validity and missingness. It
then reports signal fidelity, dynamic fidelity, behavioral fidelity, learning and
transfer, and resources. The report next gives ablation contrasts and robustness under
noise, prompt or stimulus perturbation, and alternative encodings. Only after these
sections does it discuss a possible quantum mechanism. This order prevents an attractive
quantum interpretation from hiding a basic failure of data or task validity.

The protocol explicitly supports negative outcomes. If no condition improves over C0,
the result may indicate that the target is well handled classically. If Q1 improves over
C0 but not C1, the result is a QML-versus-classical observation with no advantage over a
brain-inspired spiking model. If Q1 improves in ideal simulation but fails in Q2, the
result is a resource and noise boundary. If Q1 and A2 are indistinguishable, the
quantum feature map may not be necessary. If Q1 matches behavior but not dynamics, the
benchmark reports behavioral imitation without mechanistic imitation.

This protocol is also designed for later replication. A replication may change the
hardware or dataset, but it must preserve the task card, split logic, primary metrics,
and contrast definitions. Changes in representation are reported as a new benchmark
version. Scores from different versions are not pooled unless the mapping of tasks and
representations is documented. A versioned contamination ledger records later releases
of neural data or circuit libraries that could have entered model training.

\section{Interpretation, Safety, and Limits}
\label{sec:limits}

The decision matrix has three levels. Support for useful quantum assistance requires a
Q1 improvement over both C0 and C1 on a declared metric under a matched resource budget.
Support for a quantum-resource mechanism additionally requires degradation in A1 or A2
and persistence under Q2. Support for brain-like imitation requires a multi-axis match
including perturbation response and held-out transfer. A model can satisfy none, one, or
some of these levels.

The benchmark cannot establish quantum consciousness. A task model may have a quantum
internal representation while lacking experience; a classical model may match a neural
trajectory while lacking biological implementation. Functional fidelity is therefore
reported with its target scale and observation model. The word ``imitate'' means only
that a declared observable or behavior is reproduced within a declared tolerance.

Safety requires offline or sandboxed execution, no medical or identity inference, no
unreviewed scientific deployment, and no claims about a person's mental state. Neural
data are de-identified and governed by consent. Quantum hardware access, if added, is
logged and rate-limited. The benchmark must not be used to rank humans or to assign
rights or responsibilities to models.

There are scientific limitations. A representation can make dynamical similarity look
better or worse. Spike binning can hide timing differences. A simulator can contain its
own modeling assumptions. Resource matching is multidimensional: equal qubit count does
not imply equal memory traffic or optimizer cost. Energy measurements vary by hardware
and queueing policy. A small benchmark can support only a local claim. None of these
limitations is fixed by adding a larger neural network or more qubits.

\section{Expected Observation Patterns}
\label{sec:patterns}

Several results would be informative without being definitive. If Q1 improves signal
prediction but not dynamics or transfer, the result supports task-specific QML utility,
not brain imitation. If Q1 matches behavior while C1 matches dynamics better, the
benchmark exposes a behavioral-mechanistic gap. If Q1 improves under ideal simulation but
loses under Q2, the result identifies a hardware and noise boundary. If Q1 and A2 are
equivalent, a classical feature map may explain the observed benefit. If Q1 degrades in
A1 but not A3, entangling operations may contribute, although implementation and
optimization confounds still require review.

For a quantum-brain probe, a positive result requires a predicted biological signature
that remains after classical stochastic and spiking alternatives are fitted. A quantum
model that merely has lower prediction error is not sufficient because flexible models
can interpolate. A negative result on one proposed observable does not disprove all
quantum-brain theories, but it narrows the tested proposal.

\begin{table}[t]
\caption{Possible profiles and bounded interpretations}
\label{tab:profiles}
\centering
\footnotesize
\begin{tabular}{p{0.37\columnwidth}p{0.49\columnwidth}}
\hline
Observed profile & Interpretation \\
\hline
Q1 high signal, low transfer & Task-specific fit; weak evidence for general imitation. \\
Q1 high dynamics, C1 lower cost & Conditional quantum utility with a practical tradeoff. \\
Q1 ideal gain, Q2 no gain & Ideal advantage does not survive hardware conditions. \\
Q1 equals A2 & Classical representation may explain the gain. \\
Q1 behavior match, dynamics mismatch & Behavioral imitation is not mechanistic imitation. \\
Quantum probe survives alternatives & Evidence for the tested observable only, not consciousness. \\
\hline
\end{tabular}
\end{table}

\section{Discussion}
\label{sec:discussion}

The value of Q-BRAIN-IMITATION-BENCHMARK is its separation of questions that are often
answered by a single demonstration. Quantum AI can be investigated as a computational
tool without making a biological claim. Brain-like AI can be evaluated at an explicit
scale without pretending to simulate every molecule. A quantum-brain hypothesis can be
treated as a source of predictions rather than as a conclusion.

The design also makes a practical recommendation. When a classical or spiking model
matches a quantum model at lower cost and equal transfer, the quantum model has no
demonstrated advantage for that target. When a quantum model wins only in an ideal
simulator, the result is a theoretical observation requiring a hardware feasibility
study. When a hybrid model helps a neural task after all controls, the appropriate claim
is conditional: quantum resources were useful for that task and resource regime.

The hardest issue is mechanistic interpretation. Neural data are observations of a
biological process, not a complete description of it. A quantum circuit may approximate
the data because quantum states are a convenient function class, not because the brain
uses the same state space. Conversely, a classical spiking model may be an excellent
functional model even if the underlying biology contains quantum effects that average
out at the measured scale. The benchmark therefore requires a bridge from model states
to observable predictions before assigning a mechanistic label.

Future versions can add longitudinal learning, embodied robots, or multimodal neural
recordings. Each extension should preserve the same principles: declare scale, match
resources, hold out concepts and interventions, log provenance, and report null results.
Any extension involving people requires ethics review and must not turn a model score
into a diagnosis of consciousness. The benchmark is a tool for disciplined comparison,
not a test for digital personhood.

\section{Conclusion}
\label{sec:conclusion}

Quantum artificial intelligence may become useful for selected neural modeling tasks,
but current evidence does not show that the human brain is a quantum computer or that a
quantum model can reproduce the brain as a whole. The proposed Q-BRAIN-IMITATION-BENCHMARK
turns the question into falsifiable comparisons among quantum, hybrid, classical, and
spiking models. It measures signal, dynamics, behavior, learning, transfer, and practical
resources separately; it includes entanglement, encoding, classical-surrogate, and
noise-aware ablations; and it requires an observable prediction before testing a
quantum-brain theory. A positive result will support a bounded computational claim. A
negative transfer or hardware result will identify a boundary. Neither outcome alone
supports a claim about consciousness or subjective experience.

\appendices
\section{Minimal Run Manifest}
\label{sec:appendix}

Each run records benchmark version, target scale, task hash, data split, model family,
model and adapter version, circuit hash, qubit count, depth, entangling-gate count,
shot count, noise model, compiler version, prompt or input hash, seed, optimizer steps,
time budget, output hash, evaluator version, and human involvement. For process claims,
the event log contains parent event, source, operation, input hash, output hash, latency,
and error category. Source is human, AI, or shared; operation is propose, revise,
accept, reject, or veto. A missing field blocks only the claim that depends on it.

The result manifest separates product tables, dynamic tables, learning tables, and
resource tables. Product tables contain signal error, behavior, and validity. Dynamic
tables contain synchrony, transition divergence, attractor occupancy, and perturbation
recovery. Learning tables contain retention, forgetting, adaptation, and held-out
transfer. Resource tables contain qubits, depth, entangling operations, shots, compile
time, run time, and optional energy. The report maps each value to source hashes and
analysis code.

The release checklist is: verify licenses; remove identifiers; freeze the data split;
publish task cards; validate circuits and hashes; run the preregistered analysis;
preserve failures; disclose unavailable weights or data; and state whether any human
selected the output. These fields make the protocol auditable while keeping its claims
strictly within functional evidence.

\begin{thebibliography}{00}
\bibitem{biamonte2017} J. Biamonte, P. Wittek, N. Pancotti, P. Rebentrost, N. Wiebe,
and S. Lloyd, ``Quantum machine learning,'' \emph{Nature}, vol. 549, pp. 195--202,
2017.
\bibitem{tegmark2000} M. Tegmark, ``Importance of quantum decoherence in brain
processes,'' \emph{Physical Review E}, vol. 61, pp. 4194--4206, 2000.
\bibitem{hameroff2014} S. Hameroff and R. Penrose, ``Consciousness in the universe:
A review of the Orch OR theory,'' \emph{Physics of Life Reviews}, vol. 11, pp. 39--78,
2014.
\bibitem{lake2017} B. M. Lake, T. D. Ullman, J. B. Tenenbaum, and S. J. Gershman,
``Building machines that learn and think like people,'' \emph{Behavioral and Brain
Sciences}, vol. 40, 2017.
\bibitem{markram2015} H. Markram \emph{et al.}, ``Reconstruction and simulation of
neocortical microcircuitry,'' \emph{Cell}, vol. 163, pp. 456--492, 2015.
\bibitem{neftci2019} E. O. Neftci, H. Mostaghimi, and A. Pedram, ``Surrogate gradient
learning in spiking neural networks,'' \emph{IEEE Signal Processing Magazine}, vol. 36,
no. 6, pp. 51--63, 2019.
\end{thebibliography}

\end{document}
